\section{\emph{Clean Architecture}: onde fica o framework}

O \citetitle{livro:martin2017} organiza-se em princípios de design
(SOLID, caps.~7 a 11), princípios de componentes (caps.~12 a 14) e
arquitetura (caps.~15 a 29). Dois capítulos importam para quem usa
Quarkus todo dia:

\begin{description}
\item[The Clean Architecture (cap.~22)] Apresenta a \emph{Dependency
  Rule}: as dependências do código-fonte apontam para dentro, das
  camadas de detalhe (web, banco, framework) para as regras de negócio,
  nunca o contrário.
\item[Frameworks Are Details (cap.~32)] O framework é um detalhe; o
  livro fala de um \enquote{casamento assimétrico} (\emph{asymmetric
  marriage}) em que você se compromete com o framework e ele não se
  compromete com você.
\end{description}

\paragraph{Ponte com Quarkus.}
O Quarkus faz muito trabalho em tempo de build: o guia
\citetitle{doc:quarkus-cdi-reference} diz que ele usa uma descoberta de
beans simplificada e tenta remover beans, interceptors e decorators não
usados durante o build. Isso funciona porque as anotações estão nas
classes que o Quarkus analisa. Uma leitura da Dependency Rule para um
projeto Quarkus: anotações de framework (\texttt{@ApplicationScoped},
\texttt{@Path}, \texttt{@Inject}) ficam nos adaptadores e na camada de
aplicação; as regras de negócio continuam sendo classes e records
comuns, testáveis com JUnit puro, como na Parte~\ref{parte:testes}. Essa divisão é
uma aplicação do livro, não uma regra do Quarkus.

O mesmo guia traz uma recomendação que só faz sentido sabendo como a JVM
acessa membros: no Quarkus, se a injeção de dependência precisa acessar
um membro \texttt{private}, tem de usar reflexão, e em executável nativo
todo membro usado por reflexão precisa ser registrado. Por isso o guia
desencoraja membros \texttt{private} em beans e sugere campos de
injeção \emph{package-private} ou injeção pelo construtor.
